\documentclass[11pt]{article}
\usepackage[margin=1.15in]{geometry}
\usepackage{lmodern}
\usepackage[T1]{fontenc}
\usepackage{amsmath,amssymb,amsthm,mathtools}
\usepackage{booktabs}
\usepackage{tabularx}
\usepackage{microtype}
\usepackage[htt]{hyphenat}
\usepackage[colorlinks=true,linkcolor=blue,citecolor=blue,urlcolor=blue]{hyperref}

\newcommand{\dd}{\,\mathrm{d}}
\newcommand{\Om}{\Omega}
\newcommand{\bd}{\partial\Omega}
\newcommand{\code}[1]{\texttt{#1}}
\newcommand{\pdn}[1]{\frac{\partial #1}{\partial n}}
\newcommand{\lammin}{\lambda_{\min}}
\newcommand{\lammax}{\lambda_{\max}}
\newcommand{\R}{\mathbb{R}}
\newcolumntype{L}{>{\raggedright\arraybackslash}X}
\theoremstyle{plain}
\newtheorem{fact}{Fact}
\newtheorem*{remark}{Remark}

\setlength{\tabcolsep}{5pt}
\renewcommand{\arraystretch}{1.15}

\title{Spectral Shape Optimization: A Mathematical Overview\\
\large How \code{douse} maximizes $\lambda_k/\lambda_1$ over polygons}
\author{douse project notes}
\date{August 2026}

\begin{document}
\maketitle

\begin{abstract}
This note is a self-contained mathematical tour of the shape-optimization problem \code{douse}
solves: maximizing the ratio $\lambda_k(\Om)/\lambda_1(\Om)$ of Dirichlet--Laplacian eigenvalues
over polygons of fixed perimeter. It states, without derivation, the Hadamard shape-derivative
formulas for simple and repeated eigenvalues; works out what these formulas imply about the
structure of optimal points of $\lambda_k/\lambda_1$, first when $\lambda_k$ is simple and then
when it sits inside a degenerate cluster; and closes by describing the optimization algorithms
built on that structure and how they are wired together in \code{douse}. For a derivation-level
treatment and implementation cross-references see \code{spectral\_optimization.tex} in the same
directory; this document is the map, that one is the territory.
\end{abstract}

\tableofcontents

\section{The problem}

Let $\Om\subset\R^2$ be bounded, open and connected, with Dirichlet spectrum
\[
  0<\lambda_1(\Om)<\lambda_2(\Om)\le\lambda_3(\Om)\le\cdots,\qquad
  -\Delta u_j=\lambda_j u_j \text{ in }\Om,\quad u_j=0\text{ on }\bd,
\]
eigenfunctions taken $L^2(\Om)$-orthonormal. We study the functional
\begin{equation}\label{eq:F}
  F(\Om)=\frac{\lambda_k(\Om)}{\lambda_1(\Om)},\qquad k\ge2\text{ fixed},
\end{equation}
over the class of $N$-gons of fixed perimeter, for a chosen $N$. Two structural facts set the
stage.

\begin{fact}[Krahn]\label{fact:krahn}
$\lambda_1$ is always simple for connected $\Om$: its eigenfunction can be taken positive, and two
orthogonal positive functions cannot coexist.
\end{fact}

\begin{fact}[Payne--P\'olya--Weinberger, proved by Ashbaugh--Benguria]\label{fact:ppw}
$\lambda_2/\lambda_1\le(j_{1,1}/j_{0,1})^2=2.5387340\ldots$ for every bounded connected planar
$\Om$, with equality only for the disk.
\end{fact}

Fact~\ref{fact:krahn} is the reason the \emph{denominator} of $F$ is always smooth: whatever
happens to $\lambda_k$, $\lambda_1$ never participates in a degeneracy. Fact~\ref{fact:ppw} bounds
$F$ itself when $k=2$ and, more importantly here, gives a certificate: any computed $F>2.5387$
is a bug, not a discovery.

Because $\Om\mapsto c\Om$ sends $\lambda_j\mapsto c^{-2}\lambda_j$, $F$ is scale invariant. The
perimeter constraint is therefore pure gauge-fixing --- it removes a redundant degree of freedom
from the parametrization without adding content to the problem. (This is unlike, say,
$\min\lambda_1$, where the same constraint \emph{is} the problem.) An $N$-gon's spectrum also
ignores translation, rotation and reflection, so \code{douse} parametrizes the shape by a vector
$p\in\R^{2N-4}$ of free vertex coordinates after fixing that gauge; the details are geometric
bookkeeping and are not needed below beyond the one fact that $p\mapsto\Om(p)$ is a smooth,
locally invertible map on an open (but not all of $\R^{2N-4}$) admissible set.

\section{The Hadamard formulas}\label{sec:hadamard}

Everything that follows rests on how $\lambda_j$ responds to a perturbation of the domain. Let
$V$ be a velocity field on $\bd$ (only its normal component $V\cdot n$ matters) and let $\Om_t$
be the domain swept out by transporting $\bd$ along $V$.

\begin{fact}[Hadamard, simple eigenvalue]\label{fact:had-simple}
If $\lambda$ is a simple eigenvalue with $\|u\|_{L^2(\Om)}=1$, then $t\mapsto\lambda(\Om_t)$ is
differentiable at $t=0$ and
\begin{equation}\label{eq:hadamard-simple}
  \frac{\dd\lambda}{\dd t}\bigg|_{t=0}=-\int_{\bd}\Big(\pdn{u}\Big)^{\!2}(V\cdot n)\dd s .
\end{equation}
\end{fact}

\begin{fact}[Hadamard, repeated eigenvalue]\label{fact:had-multi}
If $\lambda$ has multiplicity $m$ with an orthonormal eigenbasis $u_1,\dots,u_m$ for its
eigenspace, define the symmetric $m\times m$ matrix
\begin{equation}\label{eq:M}
  M(V)_{ij}=-\int_{\bd}\pdn{u_i}\pdn{u_j}(V\cdot n)\dd s .
\end{equation}
Then the $m$ eigenvalue branches bifurcating from $\lambda$ under the perturbation $t\mapsto\Om_t$
are, to first order in $t$, the eigenvalues of $M(V)$ --- i.e.\ each branch has a one-sided
derivative at $t=0$ equal to one eigenvalue of $M(V)$, the branches ordered to match.
\end{fact}

We take both facts as given; see Henrot--Pierre or Delfour--Zol\'esio for proofs. What matters for
everything below is their \emph{shape} rather than their derivation: \eqref{eq:hadamard-simple} is
a linear functional of $V\cdot n$, so a simple eigenvalue is differentiable with an honest
gradient; \eqref{eq:M} is a \emph{matrix-valued} linear functional of $V\cdot n$, so a multiple
eigenvalue is not differentiable at all in the ordinary sense --- what survives is a well-defined
first-order model for the whole cluster together, and that model is linear in $V$ even though no
individual branch is.

Fact~\ref{fact:had-simple} applied to $\lambda_1$ is unconditional, by
Fact~\ref{fact:krahn} --- $\nabla_p\lambda_1$ exists everywhere on the (open) admissible set.
Fact~\ref{fact:had-multi} applies to $\lambda_k$ wherever the shape sits at (or near) a
degeneracy, which for the regular $N$-gons that these problems tend to converge to is generic
rather than exceptional: at a regular $N$-gon, every eigenvalue past $\lambda_1$ belongs to a
cluster of size $1$ or $2$ (occasionally more, at special $N$), forced by the shape's dihedral
symmetry $D_N$.

\section{Optimal points of $\lambda_k/\lambda_1$}\label{sec:optimal}

We now ask what \eqref{eq:hadamard-simple}--\eqref{eq:M} imply about a local maximum of $F$. Write
$g_1=\nabla_p\lambda_1$, which always exists.

\subsection{$\lambda_k$ simple}\label{sec:simple}

If $\lambda_k$ is simple near $p$, $F$ is differentiable there with
\begin{equation}\label{eq:gradF}
  \nabla_p F=\frac{g_k\,\lambda_1-\lambda_k\,g_1}{\lambda_1^2},\qquad g_k=\nabla_p\lambda_k,
\end{equation}
an ordinary quotient rule. A local maximum in the interior of the admissible set is characterized,
to first order, by the single vector equation $\nabla_p F=0$, equivalently
\begin{equation}\label{eq:crit-simple}
  \lambda_1\,g_k=\lambda_k\,g_1 ,
\end{equation}
i.e.\ the two shape gradients are parallel with a ratio fixed by the current value of $F$ itself.
Geometrically: at a critical point of $F$, the boundary velocity field that most increases
$\lambda_k$ is (up to sign and scale) the same field that most increases $\lambda_1$ --- there is
no direction that helps one without proportionally helping the other. This is a codimension-$1$-worth-of-conditions
statement in each of the $2N-4$ parameter directions and is otherwise unremarkable: it is exactly
what one gets by differentiating any smooth ratio, with the geometric content entirely loaded into
what $g_1$ and $g_k$ \emph{are}. Second-order conditions (a negative-semidefinite Hessian of $F$
restricted to the admissible tangent space) then distinguish maxima from saddle points in the
usual way, and are not special to this problem.

The interesting mathematics starts precisely where \eqref{eq:crit-simple} stops applying: at the
shapes these optimizations actually converge to, $\lambda_k$ is essentially never simple.

\subsection{$\lambda_k$ degenerate: the directional derivative}\label{sec:degen-general}

Suppose $\lambda_k$ sits in a cluster of multiplicity $m$, occupying position $j\in\{0,\dots,m-1\}$
when the cluster's members are listed in increasing order (so $j=0$ is the bottom of the cluster,
$j=m-1$ the top). By Fact~\ref{fact:had-multi}, the one-sided directional derivative of
$\lambda_k$ in a parameter direction $d$ is $\sigma_j(M(d))$, the $j$-th ordered eigenvalue of the
affine matrix function $M(d)=\sum_i d_i M^{(i)}$ built from \eqref{eq:M} applied to the $2N-4$
coordinate velocity fields. Combining with the (always-defined) derivative of $\lambda_1$, the
one-sided directional derivative of $F$ itself is
\begin{equation}\label{eq:dF}
  \dd F(d)=\frac{\sigma_j\big(M(d)\big)\,\lambda_1-\lambda_k\,(g_1\cdot d)}{\lambda_1^2}.
\end{equation}
(When the cluster's members are not exactly equal but merely close, the correct model replaces
$M(d)$ by $\Lambda+M(d)$ with $\Lambda=\mathrm{diag}(\lambda_\ell-\lambda_k,\dots)$ recording the
actual splitting; this refines \eqref{eq:dF} without changing anything said below, and is taken up
in \S\ref{sec:algorithms}.)

Equation~\eqref{eq:dF} is positively homogeneous of degree $1$ in $d$ (since $\sigma_j$ is), so
$F$ is directionally, but not (in general) Gateaux, differentiable at $p$: $\dd F(-d)\ne-\dd F(d)$
whenever $\sigma_j$ is genuinely nonlinear in $d$, which happens exactly when $0<m\le n_p$ and $j$
is not both the top and bottom of a singleton cluster. $F$ has a \emph{kink} at every such shape,
the multivariate analogue of $|x|$ at $x=0$: it is the norm of the shape gradient near the corner
$V=0$ of the eigenvalue branches that produces the kink, not any failure of the underlying spectral
data to vary smoothly.

\subsection{First-order optimality at a degeneracy}\label{sec:optimality-degen}

At a local maximum, no direction can increase $F$ to first order:
\begin{equation}\label{eq:opt-degen}
  \dd F(d)\le0\qquad\text{for every admissible direction }d .
\end{equation}
Because $\sigma_j$ is positively homogeneous, \eqref{eq:opt-degen} is equivalent to its restriction
to $\|d\|=1$, and it is exactly this restricted problem --- maximize $\dd F$ over a sphere (or, with
active linear constraints, a sphere cap) --- whose \emph{value} is what a numerical certificate has
to test: $p$ is first-order optimal iff that maximum is $\le0$. The structure of that maximization,
and hence of the certificate and of the ascent step that precedes it, depends entirely on where
$\lambda_k$ sits inside its cluster:

\begin{center}
\begin{tabularx}{\textwidth}{@{}lLL@{}}
\toprule
position $j$ & $\sigma_j(M(d))$ as a function of $d$ & consequence for \eqref{eq:opt-degen}\\
\midrule
$j=0$ (bottom of cluster) & concave & concave maximization: a genuine, tractable optimum\\
$j=m-1$ (top of cluster) & convex & maximizing a convex function over a ball: the max sits on the
  boundary and the feasible set of \emph{improving} directions is generically large unless a
  further reduction applies (\S\ref{sec:algorithms})\\
$0<j<m-1$ (interior) & neither & no closed form; genuinely nonconvex, though cheap to evaluate\\
$m=1$ (simple) & linear & recovers \S\ref{sec:simple}\\
\bottomrule
\end{tabularx}
\end{center}

This dichotomy is the load-bearing fact of the whole theory: $\lammin$ of a symmetric affine
matrix function is concave and $\lammax$ is convex, a classical and easy fact (each is an infimum,
respectively supremum, of linear functions $u^\top M(d)u$ over unit vectors $u$), and it is what
makes the bottom of a cluster qualitatively different from the top.

\paragraph{Which case actually occurs.} At the regular $N$-gons that these optimizations converge
to, the pattern (forced by $D_N$ symmetry and confirmed numerically) is $\lambda_1$ simple, then
clusters of size $1$ or $2$ alternating, with $\lambda_2$ always the \emph{bottom} of its pair.
Maximizing $\lambda_2/\lambda_1$ therefore lands in the tractable concave row for every $N$; the
higher ratios $\lambda_4/\lambda_1,\dots$ can land in any row depending on $N$ and $k$, which is
why \code{douse} needs all three branches of \S\ref{sec:algorithms} rather than only the concave
one.

\subsection{A worked degeneracy: when the first-order test is vacuous}\label{sec:vacuous}

It is worth seeing explicitly that \eqref{eq:opt-degen} can hold identically, carrying no
information, which is a genuine phenomenon and not a numerical artifact. Take $k=2$, $m=2$ at a
regular $N$-gon. Two facts hold there exactly:
\begin{enumerate}
\item $g_1=0$ --- the regular $N$-gon minimizes $\lambda_1$ at fixed perimeter among $N$-gons
  (Polya--Szeg\H{o}, $N=3,4$; conjectured/observed more generally), so the denominator's gradient
  vanishes at exactly the shapes where the numerator degenerates;
\item $\operatorname{tr}M(d)=0$ for every $d$ --- the symmetric function $\lambda_2+\lambda_3$ is
  stationary there, a consequence of \eqref{eq:M} and the shape's symmetry.
\end{enumerate}
With $g_1=0$, \eqref{eq:dF} collapses to $\dd F(d)=\sigma_0(M(d))/\lambda_1=\lammin(M(d))/\lambda_1$.
Because $M(-d)=-M(d)$, the sign condition $\lammin(M(d))\le0$ for all $d$ is equivalent to $M(d)$
never being positive definite, and for a $2\times2$ block exactly to $\det M(d)\le0$ for all $d$.
The traceless condition forces $M(d)=\begin{psmallmatrix}\alpha&\beta\\\beta&-\alpha\end{psmallmatrix}$
with $\alpha,\beta$ linear in $d$, so $\det M(d)=-\alpha^2-\beta^2\le0$ \emph{identically} ---
\eqref{eq:opt-degen} is satisfied automatically, for every regular $N$-gon, and certifies nothing
about optimality: the quadratic form behind it is negative semidefinite of rank at most $2$ no
matter how many free parameters $n_p$ the shape has. The verdict is then decided at second order,
on the null space where the first-order test is silent; only for $N=3$ does that null space vanish
(there $n_p=2$ exactly matches the rank), making the equilateral triangle a genuine first-order
\emph{strict} local maximizer of $\lambda_2/\lambda_1$ among triangles.

The moral for the rest of the document: a first-order stopping test at a cluster is a real
mathematical statement, not a formality, and it can degenerate into a tautology at exactly the
symmetric points the optimization is trying to find. \code{douse} treats this as a fact to be
checked (\code{certify.py}), not something the ascent algorithm itself can be trusted to notice.

\subsection{Kinks that are geometric, not spectral}\label{sec:kink}

One further nonsmooth phenomenon deserves separating out because it looks like the eigenvalue
degeneracy of \S\ref{sec:degen-general} but has nothing to do with it: $F$ also has kinks at
shapes where the polygon itself degenerates, e.g.\ where a vertex becomes flat (an $N$-gon
momentarily coinciding with an $(N-1)$-gon). There $\lambda_k$ is simple and smooth in the shape,
but the exponent $\nu=\pi/\omega$ governing the corner behavior of the eigenfunction crosses $1$
as the interior angle $\omega$ passes through $\pi$, and the two sides (reentrant, $\nu<1$;
convex, $\nu>1$) are different analytic problems meeting at a point. The result is again a
concave kink with finite, unequal, nonzero one-sided derivatives, and again a first-order model
cannot land on it. This is a fact about the \emph{parametrization} intersecting a fact about
corner singularities, not about the multiplicity structure of \S\ref{sec:degen-general}, and
\code{douse} handles the two with genuinely different machinery (\S\ref{sec:algorithms}).

\section{Optimization algorithms}\label{sec:algorithms}

Sections~\ref{sec:hadamard}--\ref{sec:optimal} describe the local model available at a point $p$:
a gradient when $\lambda_k$ is simple, a matrix pencil when it is not, always with $g_1$ available.
\code{douse} turns that model into an ascent method by a standard device --- trust-region maximization
of the local model, followed by an accept/reject test against the true objective --- specialized in
three ways to the structure above.

\subsection{Trust-region ascent, in outline}

At iterate $p$ with radius $\Delta$, propose $d^\star=\arg\max_{\|d\|\le\Delta}(\text{local model
of }\dd F)$, evaluate $F(p+d^\star)$ exactly (one eigensolve), and accept or shrink according to
the usual ratio test $\rho=(F(p+d^\star)-F(p))/(\text{predicted gain})$. Because the local model is
positively homogeneous of degree $1$ (\S\ref{sec:degen-general}), its optimal value over
$\|d\|\le\Delta$ scales linearly in $\Delta$; the stopping test is therefore not the raw optimal
value but \emph{criticality} $=(\text{optimal gain})/\Delta$, the best achievable directional
derivative, which is the quantity \eqref{eq:opt-degen} actually tests and which is scale-free.

\subsection{The smooth phase: BFGS/SR1 inside an exact trust-region solve}

Away from any degeneracy (\S\ref{sec:simple}), the local model is linear and its maximizer over a
ball is trivial ($d\propto g_k\text{ or }\nabla F$), which wastes the curvature information a run
accumulates for free along the way. \code{douse} instead fits quadratic models
$g\cdot d+\tfrac12 d^\top Bd$ to $-\lambda_1$ and to the (possibly still-active) cluster mean, using
symmetric-rank-one updates (SR1, chosen over BFGS because neither scalar is known to be convex
along the ascent path) built from consecutive \emph{accepted} smooth gradients, and solves the
resulting trust-region subproblem exactly via the More--Sorensen secular-equation characterization
$d(\sigma)=-(B+\sigma I)^{-1}g$ --- cheap here since the parameter dimension is at most a few dozen.
The curvature chain is deliberately broken (not discarded) whenever the model switches into cluster
mode, since a secant pair straddling a nonsmooth excursion compares two things that are not the
same function.

\subsection{The cluster phase: matching the subproblem to $j$}

When $\lambda_k$ is (nearly) degenerate, \S\ref{sec:optimality-degen}'s three cases become three
different subproblem solvers, all maximizing \eqref{eq:dF} over $\|d\|\le\Delta$:

\begin{itemize}
\item \textbf{Bottom of cluster ($j=0$), concave.} For a $2\times2$ block the matrix inequality
  defining $\lammin\ge t$ reduces to two scalar conditions on trace and determinant, solved as a
  small smooth program; equivalently (and more robustly, avoiding a degenerate constraint gradient
  at the natural starting point) $\lammin$ of a $2\times2$ symmetric matrix is written directly as
  \emph{linear minus a norm}, turning the subproblem into an unconstrained-looking concave program
  in $d$ alone. For $m\ge3$ blocks the honest semidefinite program is solved numerically.
\item \textbf{Top of cluster ($j=m-1$), convex.} Maximizing a convex function over a ball puts the
  optimum on the sphere, but for $m=2$ the identity
  $\lammax=\ell\cdot d+\|Ld\|$ (with $\ell,L$ linear in $d$) lets the maximization over
  $\|d\|=\Delta\subset\R^{n_p}$ be exchanged for a maximization over the unit circle in $\R^2$: a
  one-dimensional search, exact. For $m\ge3$ the same reduction becomes a maximization over the
  sphere $S^{m-1}$, handled by multi-start.
\item \textbf{Interior of cluster ($0<j<m-1$).} Neither reduction applies, and the subproblem is
  genuinely nonconvex; but since evaluating the model at any trial $d$ is a cheap $m\times m$
  eigenvalue problem, it is simply maximized directly by multi-start local optimization, seeded
  from the two extreme-case solutions above. This branch carries no first-order certificate of its
  own, only empirical agreement with brute-force search.
\end{itemize}

\noindent Whichever branch produces the step, the criticality test of \S\ref{sec:optimality-degen}
is always evaluated against the exact model \eqref{eq:dF}, so a heuristic interior step can still
be checked against the certificate it does not itself supply.

Near, but not exactly at, a degeneracy the cluster is only \emph{$\varepsilon$-active}: its
members are close, not equal, and the correct affine model is $\Lambda+M(d)$ rather than $M(d)$
alone, with $\Lambda$ the diagonal matrix of the actual (small) splitting. Treating a near-degeneracy
as an exact one makes the existing splitting look free and produces a systematically over-optimistic
predicted gain; the width of the active set is tied to the trust radius so that it neither
under-clusters (which manufactures a fictitious O(1) gradient at a genuinely nonsmooth point) nor
over-clusters indefinitely as the run converges.

\subsection{Second-order refinement of the ratio itself}

The model \eqref{eq:dF} linearizes numerator, denominator, and quotient all at once, which is more
loss than necessary: the quotient's own curvature is exact arithmetic once the two pieces are in
hand, so \code{douse} additionally forms
\[
  \widehat F(d)=\frac{\lambda_k+\sigma_j(\Lambda+M(d))+\tfrac12 d^\top H_\mu d}
                     {\lambda_1+g_1\cdot d+\tfrac12 d^\top H_1 d}-\frac{\lambda_k}{\lambda_1},
\]
carrying SR1 curvature $H_1,H_\mu$ on the two genuinely smooth scalars --- $\lambda_1$ and the
cluster mean $\mu$ --- while leaving the nonsmooth eigenvalue term $\sigma_j$ at its exact,
first-order-only, form (a second-order model of a kink is fiction). The linear/matrix-pencil steps
above are retained as seeds for maximizing $\widehat F$, so the refinement can only improve on them,
never lose ground.

\subsection{Crossing geometric kinks}

The kinks of \S\ref{sec:kink} are not reachable by any local step, however the trust region is
sized, because they sit at a discontinuity in the parametrization itself (a flat vertex, or a
collapsed edge, correspond to different regions of $p$-space that a smooth path cannot connect).
\code{douse} instead treats the two canonical degenerate configurations --- exactly flattening a
nearly-flat vertex, and deleting a very short edge --- as discrete \emph{trial moves}: propose the
candidate shape, evaluate $F$ there exactly, and accept it under the same improvement test used
for an ordinary step. This is offered only when an ordinary step has stalled at very small radius
near such a configuration, and it is what lets a run cross, for instance, from a pentagon pinned
against a minimum-edge-length constraint to the true quadrilateral optimum it is degenerating
toward.

\subsection{Assembling \texorpdfstring{$\lambda_k/\lambda_1$}{lambda\_k/lambda\_1} maximization}

Putting the pieces together, one iteration of \code{douse}'s spectral-ratio maximizer
(\code{douse.cluster\_opt.maximize\_ratio}) does the following:
\begin{enumerate}
\item Solve for the current spectrum and detect whether $\lambda_k$ is simple or belongs to an
  $\varepsilon$-active cluster, sizing the cluster width from the current trust radius.
\item Build the local model: $g_1$ always (Fact~\ref{fact:had-simple} applied to $\lambda_1$), and
  either $g_k$ (simple case) or the pencil $\{M^{(i)}\}$ and splitting $\Lambda$ (cluster case,
  Fact~\ref{fact:had-multi}).
\item Dispatch to the matching trust-region subproblem of \S\ref{sec:algorithms}, refine with the
  second-order quotient model, and test the resulting step's criticality against the linear model;
  stop (and report an achieved \emph{resolution}, not a bare pass/fail) if criticality is at or
  below tolerance, or if it has stopped changing as the radius shrinks --- a signal that the floor
  has been reached is the accuracy of the eigenvalue solves themselves, not that the point is a
  first-order maximizer to arbitrary precision.
\item Otherwise take the step (subject to any geometric constraints, linearized into the same
  subproblem), evaluate the true ratio, and accept or shrink by the ratio test; if the step is
  instead rejected while the shape is pinned near a geometric kink, offer the discrete crossing
  moves of \S\ref{sec:kink}.
\item On termination, re-solve the eigenproblem once from a cold start at the reported optimum,
  since a continuation-tracked spectrum converging \emph{onto} a degeneracy is exactly where
  tracking is least trustworthy, and report the invariants of \S\ref{fact:krahn}--\ref{fact:ppw}
  as a sanity check on the answer.
\end{enumerate}

The upshot is that the algorithm's control flow mirrors the mathematics of
\S\ref{sec:optimal} directly: every branch it can take corresponds to one row of the
position-in-cluster table, and the two places it departs from a plain smooth trust-region method
--- the matrix-pencil subproblems and the discrete kink-crossing moves --- are exactly the two
places \S\ref{sec:optimal} identified $\lambda_k/\lambda_1$ as failing to be an ordinary smooth
function of the shape.

\section*{Pointers}

\noindent\begin{tabularx}{\textwidth}{@{}lL@{}}
\toprule
\S\ref{sec:hadamard} & \code{douse/derivative.py}; the underlying quadrature is \code{lappy}'s
  \code{hadamard\_quad} \\
\S\ref{sec:degen-general}--\S\ref{sec:vacuous} & \code{douse/derivative.py}
  (\code{cluster\_matrices}), \code{douse/certify.py} \\
\S\ref{sec:kink} & \code{douse/degeneracy.py} \\
\S\ref{sec:algorithms} & \code{douse/cluster\_opt.py} (\code{local\_model},
  \code{trust\_region\_step}, \code{maximize\_ratio}), \code{douse/quasi\_newton.py} \\
\bottomrule
\end{tabularx}

\end{document}
